\documentclass[../main.tex]{subfiles}
\graphicspath{{\subfix{../img/}}}
\begin{document}
	\subsection{Introduzione}
	L'agire razionale rispetto alla scopo è un agire orientato sulla base della valutazione dei mezzi più adeguati per raggiungere un certo scopo. \\
	Agente razionale usa l'intelligenza come capacità di interagire con un ambiente.\\ \\
	Caratteristiche degli agenti:
	\begin{itemize}
		\item Sono situati (contextualizzati):
		\begin{itemize}
			\item Ricevono percezioni da un ambiente
			\item Agiscono sull'ambiente mediante azioni (attuatori)
		\end{itemize}
		\item Hanno abilità sociale:
		\begin{itemize}
			\item Capaci di comunicare
		\end{itemize}
		\item Hanno credenze, obiettivi, intenzioni
		\item Sono embodied:
		\begin{itemize}
			\item Sono limitati da un corpo e potrebbero essere soggetti ad emozioni
		\end{itemize}
	\end{itemize}
	Un agente percepisce il proprio ambiente attraverso sensori ed agisce attraverso attuatori.
	\begin{figure}[h]
		\centering
		\adjustbox{cfbox=white 5pt 0cm}{\includegraphics[scale=0.3]{6.jpeg}}
		\caption{Schema astratto funzionamento agente}
		\label{fig:im-6}
	\end{figure}\\ \\
	Terminologia agenti:
	\begin{itemize}
		\item \textbf{Percezione}: input da sensori.
		\item \textbf{Sequenza percettiva}: Storia completa delle percezioni.
		\item \textbf{Funzione agente}:
		\begin{itemize}
			\item Definisce l'azione da compiere per ogni sequenza percettiva (descrive completamente l'agente);
			\item Implementa da un programma agente.
		\end{itemize}
		\item \textbf{Compito (IA)}: progettare il programma agente.
	\end{itemize}
	Serve un criterio di valutazione oggettivo dell’effetto delle azioni dell’agente (della sequenza di stati dell’ambiente).\\ \\
	Misura di prestazione:
	\begin{itemize}
		\item Esterna (come vogliamo che il mondo evolva?);
		\item Scelta dal progettista a seconda del problema considerando l’effetto desiderato sull'ambiente;
		\item (Possibile) Valutazione su ambienti diversi.
	\end{itemize}
	\subsubsection{Esempio}
	Fissata la misura delle performance possiamo valutare come opera nel tempo sull’ambiente.
	\begin{figure}[h]
		\centering
		\adjustbox{cfbox=white 5pt 0cm}{\includegraphics[scale=0.3]{7.jpeg}}
		\caption{Schema esempio}
		\label{fig:im-7}
	\end{figure}
	Così possiamo dire che la razionalità è relativa a:
	\begin{itemize}
		\item La misura di prestazioni;
		\item Le conoscenze pregresse dell’ambiente;
		\item Le percezioni presenti e passate (seq. percettiva);
		\item Le capacità dell’agente (azioni possibili).
	\end{itemize}
	\subsection{Agenti Intelligenti/razionale}
	\textbf{Agente razionale}: Per ogni sequenza di percezioni compie l’azione che massimizza il valore atteso della misura delle prestazioni, considerando le sue percezioni passate e la sua conoscenza pregressa.\\ \\
	Razionalità non onniscienza:
	\begin{itemize}
		\item Non si pretendono perfezione e conoscenza del “futuro”, ma massimizzare il risultato atteso
		\item Potrebbero essere necessarie azioni di acquisizione di informazioni o esplorative
	\end{itemize}
	Razionalità non onnipotenza:
	\begin{itemize}
		\item Le capacità dell’agente sono limitate
	\end{itemize}
	Razionalità e apprendimento:
	\begin{itemize}
		\item Raramente tutta la conoscenza sull’ambiente può essere fornita “a priori” (dal programmatore)
		\item L’agente razionale deve essere in grado di modificare il proprio comportamento con l’esperienza (le percezioni passate)
		\item Può migliorare esplorando, apprendendo, aumentando autonomia per operare in ambienti differenti o mutevoli
	\end{itemize}
	\subsection{Agenti e ambiente}
	\textbf{Ambiente}: Dove l’agente opera, tutto ciò serve a definire un problema per un agente (Agente razionale = soluzione).\\ \\
	\textbf{PEAS}:
	\begin{itemize}
		\item P: Performance/prestazione;
		\item E: Enviroment/ambiente;
		\item A: Actuators/attuatori;
		\item S: Sensors/sensori.
	\end{itemize}
	\subsubsection{Esempi}
	\begin{tabular}{|>{\raggedright\arraybackslash}p{3cm}|>{\raggedright\arraybackslash}p{3cm}|>{\raggedright\arraybackslash}p{3cm}|>{\raggedright\arraybackslash}p{3cm}|>{\raggedright\arraybackslash}p{3.5cm}|}
		\hline
		\textbf{Problema} & \textbf{P} & \textbf{E} & \textbf{A} & \textbf{S} \\
		\hline
		Diagnosi medica & Diagnosi corretta & Pazienti, ospedale & Domande, suggerimenti test, diagnosi & Sintomi, Test clinici, risposte paziente \\
		\hline
		Analisi immagini & \% img/zone correttamente classificate & Collezione di fotografie & Etichettatore di zone nell’immagine & Array di pixel \\
		\hline
		Robot “selezionatore” & \% delle parti correttamente classificate & Nastro trasportatore & Raccogliere le parti e metterle nei cestini & Telecamera (pixel di varia intensità) \\
		\hline
		Giocatore di calcio & Fare più goal dell’avversario & Altri giocatori, campo di calcio, porte & Dare calci al pallone, correre & Locazione pallone altri giocatori, porte \\
		\hline
	\end{tabular}
	\subsection{Proprietà dell’ambiente-problema}
	\begin{itemize}
		\item Completamente/parzialmente osservabile
		\item Agente singolo/multi-agente
		\item Deterministico/stocastico/non deterministico
		\item Episodico/sequenziale
		\item Statico/dinamico
		\item Discreto/continuo
	\end{itemize}
	\textbf{Ambiente completamente osservabile}:
	\begin{itemize}
		\item L’apparato percettivo è in grado di dare una conoscenza completa dell’ambiente o almeno tutto quello che serve a decidere l’azione;
		\item Non c’è bisogno di mantenere uno stato del mondo (esterno).
	\end{itemize}
	\textbf{Ambiente parzialmente osservabile}:
	\begin{itemize}
		\item Sono presenti limiti o inaccuratezze dell’apparato sensoriale.
	\end{itemize}
	\textbf{Ambiente singolo/multiagente}:
	\begin{itemize}
		\item \textbf{Distinzione agente/non agente}: Il mondo può anche cambiare per eventi, non necessariamente per azioni di agenti.
		\item \textbf{Ambiente multi-agente competitivo (schacchi)}: Comportamento randomizzato (è razionale);
		\item \textbf{Ambiente multi-agente cooperativo (o benigno)}: Stesso obiettivo Comunicazione.
	\end{itemize}
	\textbf{Predicibilità dell’ambiente:}
	\begin{itemize}
		\item \textbf{Deterministico}: Se lo stato successivo è completamente determinato dallo stato corrente e dall’azione. Esempio: scacchi;
		\item \textbf{Stocastico}: Esistono elementi di incertezza con associata probabilità. Esempi: guida, tiro in porta;
		\item \textbf{Non deterministico}: Si tiene traccia di più stati possibili risultato dell’azione (ma non in base ad una probabilità).
	\end{itemize}
	\textbf{Ambiente statico/dinamico:}
	\begin{itemize}
		\item \textbf{Statico}: Il mondo non cambia mentre l’agente decide l’azione
		\item \textbf{Dinamico}: Cambia nel tempo, va osservata la contingenza. Tardare equivale a non agire
		\item \textbf{Semi-dinamico}: L’ambiente non cambia ma la valutazione dell’agente si. Esempio: Scacchi con timer
	\end{itemize}
	\textbf{Discreto/continuo}\\
	Possono assumere valori discreti o continui:
	\begin{itemize}
		\item Lo stato: solo un numero finito di stati
		\item Il tempo
		\item Le percezioni
		\item Le azioni
	\end{itemize}
	\textbf{Noto/ignoto:}
	\begin{itemize}
		\item Distinzione riferita allo stato di conoscenza dell’agente
		\item L’agente conosce l’ambiente oppure deve compiere azioni esplorative?
		\item Noto diverso da osservabile (e.g. carte coperte, regole note)
	\end{itemize}
	\textbf{Ambienti reali}: Parzialmente osservabili, stocastici, sequenziali, dinamici, continui, multi-agente, ignoti.
	\subsection{Struttura}
	Quattro tipi fondamentali in ordine di generalità crescente:
	\begin{itemize}
		\item Semplici agenti riflessi
		\item Agenti riflessi con stato
		\item Agenti basati su obiettivi
		\item Agenti basati sull'utilità
	\end{itemize}
	Tutti questi possono essere trasformati in agenti con apprendimento (learning agents).\\
	Agente = Architettura + Programma
	\[Ag: \; \underset{percezioni}{P} \rightarrow \underset{azioni}{Az}\]
	\begin{center}
		Il programma dell'agente implementa la funzione $Ag$\\
		$\downarrow$\\
		\textbf{function} Skeleton-Agent (\textit{percept}) \textbf{returns} action\\
		\textbf{static}: memory, the agent’s memory of the world
		\begin{align*}
			\textit{memory} &\leftarrow \textit{UpdateMemory}(\textit{memory}, \textit{percept}) \\
			\textit{action} &\leftarrow \textit{Choose-Best-Action}(\textit{memory}) \\
			\textit{memory} &\leftarrow \textit{UpdateMemory}(\textit{memory}, \textit{action})\\
			\textbf{return} & \textit{ action}
		\end{align*}
	\end{center}
	\textbf{Agente basato su tabella}: La scelta dell’azione è un accesso a una tabella che associa un’azione ad ogni possibile sequenza di percezioni.\\
	Ha diversi problemi come: la dimensione, era difficile da costruire, no autonomia, difficile aggiornamento.\\ \\
	\textbf{Agenti reattivi semplici:}
	\begin{itemize}
		\item Questi agenti selezionano le azioni sulla base della percezione corrente, ignorando il resto della cronologia della percezione;
		\item Gli agenti reattivi hanno la proprietà di essere semplici, ma sono anche di intelligenza limitata;
		\item L’agente funzionerà solo se la decisione corretta può essere presa solo sulla base della percezione corrente, cioè solo se l’ambiente è completamente osservabile. Anche un minimo di inosservabilità può causare seri problemi.
	\end{itemize}
	\textbf{Agenti con stato interno (modello)}: Indipendentemente dal tipo di rappresentazione utilizzata, è raramente possibile che l’agente determini esattamente lo stato attuale di un ambiente parzialmente osservabile.\\
	La casella etichettata "com’è il mondo adesso" rappresenta la "migliore ipotesi" dell'agente (o talvolta la migliore ipotesi).\\
	\textbf{Agenti con obiettivo:}
	\begin{itemize}
		\item Ove bisogna pianificare una sequenza di azioni per raggiungere l'obiettivo (goal)
		\item A volte la selezione dell'azione basata sull'obiettivo è semplice, es. quando si raggiunge l'obiettivo con una singola azione
		\item Altre volte sarà più complicato, ad esempio, quando l'agente deve considerare lunghe sequenze di azioni per trovare un modo per raggiungere l'obiettivo
		\item Un processo decisionale di questo tipo è fondamentalmente diverso dalle regole condizione-azione, in quanto implica le considerazioni del futuro {\color{red} (Da approfondire!! [secondo pacco di slide da pagina 47])}
		\item Sono guidati da un obiettivo nella scelta dell’azione
		\item Devono pianificare una sequenza di azioni per raggiungere l’obiettivo
		\item Meno efficienti ma più flessibili di un agente reattivo
	\end{itemize}
	\textbf{Agenti con funzione di utilità:}
	\begin{itemize}
		\item La funzione di utilità è essenzialmente un'interiorizzazione della misura della performance
		\item Se la funzione di utilità interna e la misura della prestazione esterna sono in accordo, allora un agente che sceglie le azioni per massimizzare la propria utilità sarà razionale in base alla misura della prestazione esterna
		\item Un agente basato sull'utilità deve modellare e tenere traccia del suo ambiente, compiti che hanno comportato una grande quantità di ricerca sulla percezione, rappresentazione, ragionamento e apprendimento.
	\end{itemize}
	\textbf{Agenti che apprendono (learning agent):}
	\begin{enumerate}
		\item Componente di apprendimento:
		\begin{enumerate}
			\item Produce cambiamenti al programma agente; 
			\item Migliora prestazioni, adattando i suoi componenti, aprendo dall'ambiente;
		\end{enumerate}
		\item Elemento esecutivo (performance element):
		\begin{enumerate}
			\item Il programma agente;
			\item Responsabile del selezionare le azioni;
		\end{enumerate}
		\item Elemento critico: Osserva e dà feedback sul comportamento;
		\item Generatore di problemi: Suggerisce nuove situazioni da esplorare.
	\end{enumerate}
	\begin{figure}[h]
		\centering
		\adjustbox{cfbox=white 5pt 0cm}{\includegraphics[scale=0.3]{8.jpeg}}
		\caption{Schema Learning Agent}
		\label{fig:im-9}
	\end{figure}
	La progettazione dell'elemento di apprendimento dipende molto dalla progettazione dell'elemento di prestazione.
	\subsubsection{Esempio: taxi}
	\begin{itemize}
		\item L'elemento di prestazione consiste in una qualunque raccolta di conoscenze e procedure di cui dispone il taxi per selezionare le sue azioni di guida
		\item Il taxi esce sulla strada e guida, sfruttando questo elemento di performance
		\item Il modulo critico osserva il mondo e trasmette le informazioni all'elemento di apprendimento
		\item Da questa esperienza, l'elemento di apprendimento è in grado di formulare una regola dicendo che questa è stata una cattiva azione e l'elemento di prestazione viene modificato dall'installazione della nuova regola
		\item Il generatore di problemi potrebbe identificare alcune aree di comportamento che necessitano di miglioramenti e suggerire esperimenti, come provare i freni su superfici stradali diverse in condizioni diverse
	\end{itemize}
	\subsection{Tipi di rappresentazione dello stato agente}
	\begin{itemize}
		\item Atomica
		\item Fattorizzata (più variabili e attributi)
		\item Strutturata (aggiunge relazioni)
	\end{itemize}
	\begin{figure}[h]
		\centering
		\adjustbox{cfbox=white 5pt 0cm}{\includegraphics[scale=0.3]{9.jpeg}}
		\caption{Tipi di rappresentazione dello stato agente}
		\label{fig:im-10}
	\end{figure}
	Nella rappresentazione \textbf{atomica} ogni stato del mondo è indivisibile, non ha una struttura interna. Un singolo atomo di conoscenza è una “scatola nera” la cui unica proprietà riconoscibile è quella di essere identica o diversa da un'altra scatola nera.\\ \\
	Nella rappresentazione \textbf{fattorizzata} si suddivide ogni stato in un insieme fisso di variabili o attributi, ognuno dei quali può avere una valore. Due diversi stati fattorizzati possono condividere alcuni attributi (pos. GPS) e non altri (avere molto o non avere gas). \\
	Si passa meglio da uno stato all’altro, si può rappresentare l’incertezza.\\ \\
	Nella rappresentazione \textbf{strutturata} oggetti diversi come mucche e camion e le loro variabili descrittive e le loro relazioni possano essere descritti in modo esplicito.\\
	Quasi tutto ciò che gli esseri umani esprimono nel linguaggio naturale riguarda gli oggetti e le loro relazioni.
	\subsection{Espressività}
	Una rappresentazione più espressiva può catturare, almeno in modo altrettanto conciso, tutto ciò che può catturare una rappresentazione meno espressiva.\\
	Il linguaggio più espressivo è molto più conciso.\\
	D'altra parte, il ragionamento e l'apprendimento diventano più complessi man mano che aumenta la potenza espressiva della rappresentazione.
\end{document}